\scanpage{177}
\mspage{84}

\uline{Scholie important} Les ``pieuvres de construction'' dans le dévissage d'un espace analytique [lisse], muni d'un diviseur à croisements normaux, sont, \deleted{\unsure{homot}} à homotopie près, les espaces obtenus ainsi, fibrés sur les
\[
D_{d_*}^{*}=D_{(d_0,d_1,\ldots,d_r)}^{*}
\]
(qui sont des espaces analytiques lisses, \deleted{\unsure{qui}} associés aux parties finies \deleted{\unsure{tout}} de $\mathbb N$) et obtenus ainsi: On regarde les fibrés \deleted{\unsure{un}} $\unsure{\breve N}_{d_*}|D_{d_*}^{*}$ sur $D_{d_*}^{*}$ rel\unsure{.!} à
\[
D_{d_*}^{*}\longrightarrow D_{d_0},
\]
ce fibré \inserted{vectoriel} admet une filtration de type $d_0$, on prend un \deleted{\unsure{quotient}} gradué
\[
\unsure{\breve N}=\unsure{\breve N}_{d_*;(d'_0,d''_0)}
\]
de cette filtration \deleted{\unsure{et}} (restreint à $D_{d_*}^{*}$), et l'ouvert $\unsure{\breve N}^{*}$ --- qui est donc un fibré loc trivial sur $D_{d_*}^{*}$, à fibres $\mathbb C^{*\,d'_0-d''_0}$ --- c'est lui!

\margin{y compris la \unsure{partie vide} \unsure{correspondant} à $D_0^{*}$, mais sur $D_0^{*}$ le seul fibré à considérer est le fibré trivial isom. à $D_0$ lui-même, et de plus il n'y a pas lieu de faire de \unsure{limite} sur $D_0$ \unsure{entier}, \unsure{seulement} \unsure{aux pt} du support de $\unsure{\Theta}$.}

En se fatiguant un peu, on doit

\margin{176}
